\chapter{Weakly Interacting Gases and Inexact Differentials}

Susskind opens this lecture with a warning that sounds almost perverse: no physical device directly measures the mathematical quantity we say it measures. That is not a detachable preface. In this LazyingArt LLC and Video2Book companion chapter it is the first step of the argument, because it teaches us how to think about temperature before we ask when the ideal-gas law fails. The lecture then turns, at a student's prompting, to a dilute gas with weak interactions, and finally to a mathematical interlude on exact and inexact differentials that clarifies why heat and work are not state functions.

\section{Measurement, Calibration, and the Limited Meaning of Measuring Temperature}

No instrument reports an abstract variable by magic. A spring balance does not literally measure force; it measures the stretch of a spring. A thermometer does not literally measure temperature; it measures the height of a column. What we trust is a controlled relation between an easily observed proxy and the quantity we want.

That relation is never exact over all ranges. A spring stretched too far stops being a good force meter; at worst it breaks. Likewise, any ordinary thermometer is reliable only over a regime in which its response is smooth and approximately linear. This is the operational logic of measurement. Over a sufficiently small interval,
\begin{equation}
\delta \ell \propto \delta F,
\qquad
\delta h \propto \delta T,
\end{equation}
where \(\ell\) is the spring extension and \(h\) is the height of the thermometer column. The point is not that \(\ell\) equals \(F\), or that \(h\) equals \(T\), but that for a smooth response curve, small changes in the readout are proportional to small changes in the intended variable.

The lecture is careful about what else the readout might depend on. A spring might depend weakly on temperature, air pressure, or all sorts of nuisances that one usually ignores. A thermometer might depend weakly on gravity, orientation, or very slow changes in the material. If one of these effects matters, one folds it into the protocol: hold the thermometer vertically, use it only over a suitable range, and do not pretend that the readout is independent of every other variable in the universe.

Calibration is therefore conventional as well as physical. We choose two reference points, say the freezing and boiling points of water, and build a linear scale between them. If \(h_f\) and \(h_b\) are the corresponding column heights, then within the calibrated range one may write
\begin{equation}
T_{\mathrm{C}}=100\,\frac{h-h_f}{h_b-h_f}.
\end{equation}
This is a perfectly good temperature scale, but it already shows what Susskind wants us to notice: the zero is chosen. It is not absolute zero. A Fahrenheit thermometer reading \(0^\circ\) in winter does not mean the outside world has fallen to the bottom of thermodynamics.

\section{Equilibrium and Temperature as an Energy Scale}

The opening discussion sharpens when the class asks about equilibrium. A thermometer measures the temperature of its environment only after it has equilibrated with that environment. Before that, it is simply changing in response to it. Experimentally the rule is embarrassingly simple and completely correct: when the thermometer stops changing, it has equilibrated.

This in turn explains why the same nominal temperature can be associated with very different equilibration times. A gas may have a high temperature and an extremely low density. That does not mean anything inconsistent. It means the molecules that are present carry a large kinetic energy, but there are few of them, so they strike the thermometer only rarely. The temperature can be high while equilibration is slow. What matters for the time scale is not only temperature but also the density and thermal conductivity of the surrounding medium.

The same point appears again when the lecture turns to radiation. The radiation field in a room need not be in thermal equilibrium with the air. If matter and radiation are not equilibrated, then one must not speak as though there were a single common temperature describing both. The hidden condition behind ordinary thermodynamic language is equilibrium.

Once that condition is met, the lecture makes its conceptual pivot. Temperature should be understood not merely as a column height or a number on a scale, but as an energy scale. Later, in the ideal gas, this becomes quantitative:
\begin{equation}
E_{\mathrm{ideal}}=\frac{3}{2}NT.
\end{equation}
In the units of the lecture, \(T\) already has the dimensions of energy. At a given temperature the kinetic energy per particle is fixed, while the speed depends on the mass. That is why a heavy particle, or even a bowling ball immersed in a thermal environment, can have the same thermal energy scale as a light molecule and yet move much more slowly.

\section{From the Ideal Gas to a Weakly Interacting Gas}

The lecture's real problem begins with a student interruption: when does the ideal gas break down, and can we show it mathematically? Susskind answers in exactly the spirit of perturbation theory. We begin with a solved problem, the noninteracting gas, and then we change it a little. If the change is small enough, we expand. If the next term becomes as large as the one we started from, the approximation has broken down.

That is why the gas is taken to be weakly interacting. The molecules are dilute, typically far apart, and the range of the force is assumed short compared with the mean separation. In that regime the potential energy is small compared with the kinetic energy. We turn on the interaction with a small knob, so to speak, and then let the calculation tell us when that knob has been turned too far.

This is also why the lecture insists that the right small quantity is not introduced abstractly, but physically. The smallness comes from dilution. If the density is too high, particles are frequently near enough to feel one another and the ideal gas is poor. If the interaction reaches over very long distances, the ideal gas is poor even at low density. The perturbative expansion is justified only when both the interaction strength and its effective range are small in the appropriate sense.

With that understood, the lecture says, in effect: let us set up the problem.

\section{Pair Potentials, Pair Counting, and the Definition of \(U_0\)}

We now write the energy of a configuration as a kinetic part plus a sum over pairs. On the blackboard the notation drifts a little: \(p^2\) is written schematically inside the sum, and the mass appears once as \(M\) and later as \(m\). We normalize this in the cleaned notes and write
\begin{equation}
E=\sum_{n=1}^{N}\frac{p_n^2}{2m}+\sum_{n>m}U\!\left(|x_n-x_m|\right).
\end{equation}
Here \(x_n\) is the three-dimensional position of particle \(n\), \(p_n\) is its momentum, and \(p_n^2\) means \(p_{nx}^2+p_{ny}^2+p_{nz}^2\).

\begin{figure}[t]
\centering
\includegraphics[width=0.82\textwidth]{lecture_06_figure_04.png}
\caption{Total energy split into kinetic and pair-potential terms, with the partition function begun below.}
\end{figure}

The notation \(n>m\) is important. It excludes the fake pair \(n=m\), and it avoids double counting. The pair \((3,7)\) is the same physical pair as \((7,3)\), so only one of them belongs in the sum.

It is convenient to gather the total potential energy into a single configuration-space quantity,
\begin{equation}
U(X)=\sum_{n>m}U\!\left(|x_n-x_m|\right),
\qquad
X=(x_1,\dots,x_N),
\end{equation}
and to abbreviate the full \(3N\)-dimensional measures by
\begin{equation}
dx=\prod_{i=1}^{N}d^3x_i,
\qquad
dp=\prod_{i=1}^{N}d^3p_i.
\end{equation}
Then the partition function is
\begin{equation}
Z=\frac{1}{N!}\int dp\,dx\,
e^{-\beta\sum_{n=1}^{N}p_n^2/(2m)}\,e^{-\beta U(X)},
\qquad \beta=\frac{1}{T}.
\end{equation}
Only the kinetic exponential is clearly visible in the board frame; the extra factor \(e^{-\beta U(X)}\) is the transcript-supported completion of what the lecture is setting up.

At this point Susskind pauses over one integral and one number. The integral is the two-particle average of the pair potential over all positions:
\begin{equation}
\int d^3x_1\,d^3x_2\,U\!\left(|x_1-x_2|\right).
\end{equation}
The important number extracted from it is
\begin{equation}
U_0=\int d^3r\,U(r).
\end{equation}
The units of \(U_0\) are energy times volume. It is not just the height of the potential, and not just its range, but the combination of the two: strength times effective interaction volume.

The lecture's reconstructed relation is then
\begin{equation}
\int d^3x_1\,d^3x_2\,U\!\left(|x_1-x_2|\right)=V\,U_0.
\end{equation}
In the frame, the right-hand side is cropped after \(VU\ldots\), so \(V U_0\) should be read as the cautious clean completion rather than a literal full transcription.

\subsection{Question \& Answer}

Why does the two-particle integral scale like \(V U_0\) rather than \(V^2\)?

Because only one of the two integrations explores the whole container. If we hold one particle fixed and integrate the other one over all space, the potential is appreciable only over the small region where the two particles are close enough to feel one another. That integral is therefore not of order the whole volume; it is the interaction strength integrated over the interaction region, and that is exactly \(U_0\). Only after doing that do we move the first particle through the whole sample, which produces the single factor of \(V\).

This is also why the units work out. The full two-particle integral has units of volume squared times energy, but its large-scale dependence is only linear in the container volume; the second factor of volume is hidden inside \(U_0\).

Boundary effects are real, and the lecture insists on admitting them. A particle close to the wall cannot have a full interaction neighborhood around it. But only a small fraction of particles lie within a molecular distance of the surface, so for a macroscopic sample the correction is a surface-to-volume effect and is negligible at the level of the present approximation.

\section{Partition Function Expansion and the First Correction to the Equation of State}

Having set up the energy, the lecture announces a change of mode: now we get on autopilot. Statistical mechanics is at its most mechanical here. We write the partition function, factorize it, expand in the small quantity, and then inspect the answer afterwards to see whether it makes sense.

Because the kinetic and potential terms depend on different variables, the partition function factorizes:
\begin{align}
Z
&=\frac{1}{N!}\int dp\,dx\,
e^{-\beta\sum_n p_n^2/(2m)}\,e^{-\beta U(X)} \\
&=\left[\frac{V^N}{N!}\int dp\,e^{-\beta\sum_n p_n^2/(2m)}\right]
\left[\frac{1}{V^N}\int dx\,e^{-\beta U(X)}\right].
\end{align}
The reason for multiplying and dividing by \(V^N\) is purely practical: it isolates the old ideal-gas answer immediately. We recognize
\begin{equation}
Z_0=\frac{V^N}{N!}\int dp\,e^{-\beta\sum_n p_n^2/(2m)}
\end{equation}
as the partition function of the ideal gas, so that
\begin{equation}
Z=Z_0\left[\frac{1}{V^N}\int dx\,e^{-\beta U(X)}\right].
\end{equation}

The momentum integral is the same Gaussian integral as before:
\begin{equation}
\int dp\,e^{-\beta\sum_n p_n^2/(2m)}
=\left(\frac{2\pi m}{\beta}\right)^{3N/2}.
\end{equation}
On the board it appears as a square root raised to the power \(3N\); the cleaned form above is equivalent.

\begin{figure}[t]
\centering
\includegraphics[width=0.82\textwidth]{lecture_06_figure_05.png}
\caption{Recalled ingredients for the factorized partition function: the pair-potential integral and the Gaussian momentum factor.}
\end{figure}

The two lines visible in the figure summarize the two ingredients the lecture now reuses:
\begin{equation}
\int d^3x_1\,d^3x_2\,U\!\left(|x_1-x_2|\right)=V\,U_0,
\qquad
\left(\frac{2\pi m}{\beta}\right)^{3N/2}.
\end{equation}

The new part is the position-space factor. Since the interaction is weak, the lecture expands in the potential energy itself:
\begin{equation}
e^{-\beta U(X)}\approx 1-\beta U(X).
\end{equation}
Therefore
\begin{align}
\frac{1}{V^N}\int dx\,e^{-\beta U(X)}
&\approx \frac{1}{V^N}\int dx\,\left[1-\beta \sum_{n>m}U\!\left(|x_n-x_m|\right)\right] \\
&=1-\frac{\beta}{V^N}\sum_{n>m}\int dx\,U\!\left(|x_n-x_m|\right).
\end{align}
Now every term in the sum has the same form. Not the same value point by point, but the same integral. There is nothing special about the pair \((1,2)\) as compared with any other pair. The number of distinct pairs is
\begin{equation}
\sum_{n>m}1=\binom{N}{2}=\frac{N(N-1)}{2}\approx \frac{N^2}{2}.
\end{equation}
For one chosen pair the remaining \(N-2\) coordinate integrals give \(V^{N-2}\), and the pair integral gives \(V U_0\). Thus
\begin{align}
\frac{1}{V^N}\int dx\,e^{-\beta U(X)}
&\approx 1-\beta\,\frac{N(N-1)}{2}\,\frac{V^{N-2}(V U_0)}{V^N} \\
&=1-\beta\,\frac{N(N-1)}{2V}\,U_0 \\
&\approx 1-\beta\,\frac{N^2}{2V}\,U_0.
\end{align}
So
\begin{equation}
Z\approx Z_0\left(1-\beta\,\frac{N^2}{2V}\,U_0\right).
\end{equation}
Taking the logarithm and expanding once more for small \(U_0\),
\begin{equation}
\log Z\approx \log Z_0-\beta\,\frac{N^2}{2V}\,U_0.
\end{equation}

At this point the lecture says, in effect, that we are never interested in \(\log Z\) for its own sake. We want what we can get from it. The first choice is the energy:
\begin{equation}
E=-\frac{\partial \log Z}{\partial \beta}
\approx \frac{3}{2}NT+\frac{N^2}{2V}U_0.
\end{equation}
It is useful to factor out one power of \(N\):
\begin{equation}
E\approx N\left(\frac{3}{2}T+\frac{1}{2}\rho U_0\right),
\qquad
\rho=\frac{N}{V}.
\end{equation}
The apparently dangerous \(N^2\) has combined with the large volume to produce the density. Extensivity has been restored.

Energy is not the easiest quantity to measure in a gas, so the lecture immediately turns to pressure. Using
\begin{equation}
P=T\left(\frac{\partial \log Z}{\partial V}\right)_T,
\end{equation}
we get the ideal-gas piece from \(\log Z_0\) and the interaction correction from the new term:
\begin{equation}
P\approx \rho T+\frac{1}{2}\rho^2 U_0.
\end{equation}
The spoken sign bookkeeping becomes loose here, but the formula used in the rest of the lecture is this one: the correction is \(+\frac12 \rho^2 U_0\).

\subsection{Question \& Answer}

Are we losing anything essential by keeping only pair interactions and not triplets or higher-body effects?

We are certainly making an approximation, but for a dilute gas it is the natural first one. There may be genuine three-body forces, and even if the microscopic dynamics is built only from pair potentials, higher-order terms in the expansion probe configurations in which three or more particles happen to be close together at once. Those effects are smaller because such configurations are less probable in a dilute gas.

The lecture is cautious here, and the notes should be equally cautious. Susskind suggests that the next correction ought to involve higher powers of the density, but then stops himself and refuses to canonize a formula he has not checked. The stable point is only this: pair terms dominate the first correction, and the probability of triplets is smaller by an extra power of dilution.

\section{Interpreting the Correction and the Criterion for the Ideal Gas}

Now the lecture slows down and asks the right question: does the answer make sense? The first check is the density dependence. The ideal-gas term in the pressure is proportional to \(\rho\), because it only asks whether particles are there to hit the wall. The correction term is proportional to \(\rho^2\), because it asks for two particles to come close enough to interact. Pair physics should carry one extra power of density, and it does.

The second check is the interpretation of the energy formula. Dividing by \(N\), we have
\begin{equation}
\frac{E}{N}\approx \frac{3}{2}T+\frac{1}{2}\rho U_0.
\end{equation}
The first term is the familiar kinetic contribution per particle. The second term is an interaction energy per particle, and it depends on the density because a given particle only acquires interaction energy when another particle is nearby.

\subsection{Question \& Answer}

Why is \(\rho^2\) the right power of the density?

Because the correction is about pairs. The probability of finding one particle in a small volume is proportional to \(\rho\). The probability of finding two particles in that same small volume is proportional to \(\rho^2\). The ideal-gas pressure only needs one particle; the interaction correction needs two. That is the whole story.

The sign is equally sensible. If \(U_0>0\), the interaction is net repulsive, so compression meets extra resistance and the pressure rises:
\begin{equation}
U_0>0
\qquad \Longrightarrow \qquad
P>\rho T.
\end{equation}
If \(U_0<0\), the interaction is net attractive and the pressure is lowered:
\begin{equation}
U_0<0
\qquad \Longrightarrow \qquad
P<\rho T.
\end{equation}
This matches the lecture's qualitative discussion. Molecules are typically repulsive at very short distance and attractive in a tail farther out; one thinks here of the usual molecular potentials with a hard core and an attractive tail, the sort of profile modeled phenomenologically by a \(6\)-\(12\) potential. But statistical mechanics does not derive that profile. Molecular and atomic physics supply \(U(r)\); statistical mechanics tells us what bulk consequence follows from it.

Another point stressed in the lecture is that the pressure correction, at this order, carries no explicit temperature dependence. That makes it experimentally accessible. One can keep the gas in thermal equilibrium at fixed temperature, vary the density by changing the volume, and look for the deviation from linearity in \(P(\rho)\).

Finally we ask when the ideal gas is a good approximation. We compare the correction to the leading term:
\begin{equation}
\rho^2 U_0 \ll \rho T.
\end{equation}
Dividing through by \(\rho\),
\begin{equation}
\rho U_0 \ll T.
\end{equation}
Since \(U_0\) has units of energy times volume, the quantity \(\rho U_0\) has units of energy. It is the potential energy per particle. So the lecture's physical criterion is precise: the ideal gas works when the potential energy per particle is much smaller than the kinetic energy scale \(T\).

That is why ordinary air in a room is so well approximated by an ideal gas. One would need to compress it until molecules had a substantial probability of crowding into one another's interaction range before the correction stopped being small.

\section{Exact and Inexact Differentials: Why Heat Is Not a State Function}

At this point the lecture shifts register entirely. Susskind says, more or less, that what comes next is much less straightforward and that it confused him badly when he first learned thermodynamics. The topic is exact versus inexact differentials.

Take a function \(f(x,y)\), and think of it as the altitude of a landscape. If we take a small step \(dx,dy\), then the altitude changes by
\begin{equation}
df=\frac{\partial f}{\partial x}\,dx+\frac{\partial f}{\partial y}\,dy
=f_x\,dx+f_y\,dy.
\end{equation}
The important question is this: given two functions \(F_x(x,y)\) and \(F_y(x,y)\), when does the expression
\begin{equation}
F_x\,dx+F_y\,dy
\end{equation}
really come from the differential of a genuine scalar function?

The criterion is equality of mixed partials:
\begin{equation}
\frac{\partial F_y}{\partial x}=\frac{\partial F_x}{\partial y}.
\end{equation}
In vector language, this is the statement that the curl vanishes. In integral language, it is equivalent to path independence:
\begin{equation}
\int_{1}^{2}(F_x\,dx+F_y\,dy)
\end{equation}
depends only on the endpoints, and in particular
\begin{equation}
\oint (F_x\,dx+F_y\,dy)=0
\end{equation}
around any closed loop.

The lecture then gives two clean examples. First,
\begin{equation}
F_x=y,
\qquad
F_y=x.
\end{equation}
Here
\begin{equation}
\frac{\partial F_x}{\partial y}=1,
\qquad
\frac{\partial F_y}{\partial x}=1,
\end{equation}
so the differential is exact. Indeed,
\begin{equation}
d(xy)=y\,dx+x\,dy.
\end{equation}
Second,
\begin{equation}
F_x=y,
\qquad
F_y=-x.
\end{equation}
Now
\begin{equation}
\frac{\partial F_x}{\partial y}=1,
\qquad
\frac{\partial F_y}{\partial x}=-1,
\end{equation}
so the condition fails and the differential is inexact.

The lecture then subordinates the everyday analogies to the mathematics. Altitude is the exact case: if we walk from one point to another, the change in altitude does not depend on the path. Fuel consumption is different. One route may go straight, another over a hill, and the gas tank will not end up at the same reading. Tiredness is different for the same reason. These quantities are meaningful along a path, but they are not functions of position alone.

Now we return to thermodynamics. Here Susskind briefly says ``energy and volume'' as the independent variables, then corrects himself. The natural thermodynamic state variables for the present discussion are entropy \(S\) and volume \(V\). The first law is
\begin{equation}
dE=-P\,dV+T\,dS.
\end{equation}
With the lecture's sign convention,
\begin{equation}
dW=-P\,dV,
\qquad
dQ=T\,dS.
\end{equation}
When \(dV>0\) at positive pressure, the gas expands and does work on the piston, so the work done on the gas is negative and the gas energy decreases. That is exactly what the sign in \(dE=-P\,dV+T\,dS\) says.

To test whether heat is a state function, we rewrite
\begin{equation}
dQ=dE+P\,dV
\end{equation}
and now temporarily use \(E\) and \(V\) as coordinates in state space. Assume, for the sake of argument, that there existed a genuine state function \(Q(E,V)\). Then its derivatives would have to satisfy
\begin{equation}
\left.\frac{\partial Q}{\partial E}\right|_V=1,
\qquad
\left.\frac{\partial Q}{\partial V}\right|_E=P.
\end{equation}
If \(dQ\) were exact, the mixed partials would agree, so
\begin{equation}
\frac{\partial}{\partial V}\left(\left.\frac{\partial Q}{\partial E}\right|_V\right)
=
\frac{\partial}{\partial E}\left(\left.\frac{\partial Q}{\partial V}\right|_E\right).
\end{equation}
The left-hand side is zero, and the right-hand side is \(\left.\frac{\partial P}{\partial E}\right|_V\). Therefore exactness would require
\begin{equation}
\left.\frac{\partial P}{\partial E}\right|_V=0.
\end{equation}

Now the lecture does what it always does at the crucial moment: it checks a concrete model. For the ideal gas,
\begin{equation}
PV=NT,
\qquad
E=\frac{3}{2}NT.
\end{equation}
Eliminating \(NT\), we obtain
\begin{equation}
P=\frac{2}{3}\frac{E}{V}.
\end{equation}
Hence
\begin{equation}
\left.\frac{\partial P}{\partial E}\right|_V=\frac{2}{3V}\neq 0.
\end{equation}
The curl test fails. So \(dQ\) is not exact. Heat is not a state function. The same logic applies to work. What is exact is only \(dE\), the change in the genuine state function \(E\).

\subsection{Question \& Answer}

Why is heat not a property of the thermodynamic state?

Because heat does not tell us only where the system is; it tells us something about how the system got there. Two processes can begin at the same state and end at the same state, yet exchange different amounts of heat with the surroundings. On a closed cycle the total change in energy is zero, because the system comes back to the same state, but the total heat and the total work need not vanish separately.

That is why Susskind compares heat to fuel consumption or tiredness rather than altitude. Altitude is a property of position. Fuel consumption and tiredness record a path. Heat and work do the same. Only their combination in the first law gives a true change of state.

\section{Summary}

The lecture unfolds in two very different but tightly connected parts. First, we are reminded that measurement is indirect: a thermometer measures a proxy, and only within a controlled linear regime and after equilibration does that proxy become a reliable indicator of temperature. That prepares the later identification of temperature as an energy scale.

Then, once a student asks when the ideal gas fails, the lecture turns on a weak pair potential and lets the statistical-mechanical machinery run. The first correction to the ideal-gas law is
\begin{equation}
P\approx \rho T+\frac{1}{2}\rho^2U_0,
\end{equation}
with
\begin{equation}
\rho U_0 \ll T
\end{equation}
as the criterion for the ideal-gas approximation. The correction scales like \(\rho^2\) because it is pair physics, and its sign follows the sign of the net integrated interaction \(U_0\).

Finally, the lecture changes gears and clarifies a structural point in thermodynamics. Exact differentials belong to state functions; inexact differentials depend on path. Heat and work are inexact. Energy is exact. That distinction, once made cleanly, is not a side remark. It is part of what it means to speak precisely about thermodynamic states at all.